% ============================================================
\section{Log-Periodic Horizons (Log-Scale Coherence Towers)}
\label{sec:log-periodic}
% ============================================================

\begin{quote}\small
\textbf{What this automates.}
Many natural horizon towers are indexed by a logarithmic scale rather than a linear one: doubling resolution, octave decompositions, dyadic refinements.
This section provides the specialized framework for log-scale horizons and derives how the telescoping identity adapts to this setting.
\end{quote}

\subsection{Log-scale horizon family}

\begin{definition}[Log-scale horizon tower]
\label{def:log-horizon}\lean{LogScaleHorizonTower}
A \emph{log-scale horizon tower} on $\cH$ is a family of orthogonal projections $(\CCPi{\ell})_{\ell \in \bbR}$ (or $\ell \in \bbZ$ for the discrete case) such that:
\begin{enumerate}[label=(\roman*)]
\item \textbf{Nestedness:} $\ell \le \ell'$ implies $\CCPi{\ell} \CCPi{\ell'} = \CCPi{\ell}$.
\item \textbf{Log-scale interpretation:} The parameter $\ell$ represents a logarithmic scale, e.g., $\ell = \log_2(\text{resolution})$.
\end{enumerate}
\end{definition}

\paragraph{Model example (Fourier bandlimiting).}
\label{ex:fourier-bandlimit}
Let $\cH = L^2(\bbR)$ and define $\CCPi{\ell}$ as the Fourier bandlimiting operator to frequencies $|\xi| \le e^\ell$:
\[
(\CCPi{\ell} f)\,\widehat{}\,(\xi) :=
\begin{cases}
\hat f(\xi) & \text{if } |\xi| \le e^\ell, \\
0 & \text{otherwise}.
\end{cases}
\]
Then $\ell \le \ell'$ implies $\CCPi{\ell} \CCPi{\ell'} = \CCPi{\ell}$ (lower bandwidth is contained in higher bandwidth).

\paragraph{Model example (Dyadic resolution).}
\label{ex:dyadic}
Let $\cH = L^2([0,1])$ and define $\CCPi{n}$ (for $n \in \bbZ_{\ge 0}$) as the orthogonal projection onto piecewise-constant functions on $2^n$ equal bins.
Setting $\ell = n$ gives a discrete log-scale tower with ``resolution $= 2^\ell$.''

\subsection{Telescoping in log-scale}

\begin{theorem}[Log-scale telescoping identity]
\label{thm:log-telescoping}\lean{log_telescope_3}
Let $(\CCPi{\ell})$ be a log-scale horizon tower and let $\ell_0 < \ell_1 < \cdots < \ell_N$ be an increasing sequence of log-scale indices.
For any $x \in \cH$:
\[
\|x - \CCPi{\ell_0} x\|^2
=
\sum_{j=0}^{N-1} \|(\CCPi{\ell_{j+1}} - \CCPi{\ell_j}) x\|^2 + \|x - \CCPi{\ell_N} x\|^2.
\]
If the tower is faithful ($\CCPi{\ell} \to I$ as $\ell \to \infty$), then
\[
\|x - \CCPi{\ell_0} x\|^2 = \sum_{j=0}^{\infty} \|(\CCPi{\ell_{j+1}} - \CCPi{\ell_j}) x\|^2.
\]
\end{theorem}

\begin{proof}
This is a direct application of the Horizon Fundamental Theorem (HFT-2, Theorem~\ref{thm:HFT}) to the sequence of projections $\CCPi{\ell_j}$.
The projections are nested and orthogonal, so the increments $(\CCPi{\ell_{j+1}} - \CCPi{\ell_j}) x$ are pairwise orthogonal.
The Pythagorean theorem gives the decomposition.
\end{proof}

\begin{quote}\small
\textbf{Intuition.}
When horizons are indexed logarithmically, the telescoping identity decomposes deviation into contributions from each ``octave'' or ``scale band.''
Each term $\|(\CCPi{\ell_{j+1}} - \CCPi{\ell_j}) x\|^2$ measures how much of $x$ lives in the scale band between $\ell_j$ and $\ell_{j+1}$.
\end{quote}

\subsection{Uniform log-spacing}

\begin{corollary}[Dyadic telescoping]
\label{cor:dyadic-telescoping}\lean{DyadicTower.dyadic_telescoping}
For uniformly spaced log-indices $\ell_j = \ell_0 + j \Delta\ell$ with spacing $\Delta\ell > 0$:
\[
\|x - \CCPi{\ell_0} x\|^2 = \sum_{j=0}^{N-1} \|(\CCPi{\ell_0 + (j+1)\Delta\ell} - \CCPi{\ell_0 + j\Delta\ell}) x\|^2 + \|x - \CCPi{\ell_0 + N\Delta\ell} x\|^2.
\]
Each term corresponds to one ``log-band'' of width $\Delta\ell$.
\end{corollary}

\paragraph{Scale-invariant decomposition.}
When $\Delta\ell = \log 2$ (octave spacing), each band corresponds to a doubling of resolution.
The decomposition is then scale-invariant in the sense that each band spans the same multiplicative factor in resolution, regardless of absolute scale.

\subsection{Defect across log-scales}

The defect energy $\CCDefEnergy{\ell}{x} = \|\CCLeakExpr{d}{\ell} x\|^2$ depends on \emph{both} the shadow projection $\CCPibar{\ell}$ and the compressed state $\CCPi{\ell} x$.
As $\ell$ varies, both factors change, so monotonicity requires care.

\paragraph{Worked example (Non-monotonicity without assumptions).}
\label{ex:defect-nonmonotone}
Let $\cH = \bbR^3$, with nested tower $\CCPi{0} = e_1 e_1^*$, $\CCPi{1} = e_1 e_1^* + e_2 e_2^*$, $\CCPi{2} = I$.
Define $d = e_1 e_1^* + e_3 e_2^*$, so $d e_1 = e_1$, $d e_2 = e_3$, $d e_3 = 0$.
For $x = e_1 + e_2$:
\begin{align*}
\CCDefEnergy{0}{x} &= \|\CCPibar{0} d \CCPi{0} x\|^2 = \|\CCPibar{0} d e_1\|^2 = \|\CCPibar{0} e_1\|^2 = 0, \\
\CCDefEnergy{1}{x} &= \|\CCPibar{1} d \CCPi{1} x\|^2 = \|\CCPibar{1} d(e_1 + e_2)\|^2 = \|\CCPibar{1} (e_1 + e_3)\|^2 = \|e_3\|^2 = 1.
\end{align*}
Thus $\CCDefEnergy{0}{x} = 0 < 1 = \CCDefEnergy{1}{x}$: refining the horizon from $\CCPi{0}$ to $\CCPi{1}$ \emph{increased} the defect energy, because the newly visible component $e_2$ leaks into the shadow under $d$.

\begin{proposition}[Log-scale defect monotonicity under fixed representability]
\label{prop:log-defect}\lean{log_defect_monotonicity}
Let $d$ be a nilpotent differential ($d^2 = 0$) and $(\CCPi{\ell})$ a nested horizon tower.
Suppose $x \in \Img(\CCPi{\ell_0})$ for some $\ell_0$, i.e., $\CCPi{\ell_0} x = x$.
Then for all $\ell \ge \ell_0$:
\[
\CCDefEnergy{\ell}{x} = \|\CCPibar{\ell} d x\|^2 \quad\text{is monotone non-increasing in } \ell.
\]
\end{proposition}

\begin{proof}
Since $x \in \Img(\CCPi{\ell_0})$ and horizons are nested, $\CCPi{\ell} x = x$ for all $\ell \ge \ell_0$.
Thus $\CCDefEnergy{\ell}{x} = \|\CCPibar{\ell} d \CCPi{\ell} x\|^2 = \|\CCPibar{\ell} d x\|^2$.
For $\ell \le \ell'$, we have $\CCPibar{\ell'} = I - \CCPi{\ell'} \preceq I - \CCPi{\ell} = \CCPibar{\ell}$ in the PSD order (since $\CCPi{\ell} \preceq \CCPi{\ell'}$).
Therefore $\|\CCPibar{\ell'} d x\|^2 \le \|\CCPibar{\ell} d x\|^2$.
\end{proof}

\paragraph{Interpretation.}
The monotonicity holds when the state $x$ is ``already representable'' at some base scale $\ell_0$.
In this case, refining the horizon (increasing $\ell$) only shrinks the shadow, reducing leakage.
For states not yet representable at the base scale, $\CCPi{\ell} x$ varies with $\ell$, and defect energy may fluctuate.

\paragraph{Log-periodic structure.}
In some applications, the defect energy exhibits approximate periodicity in the log-scale variable $\ell$.
This can occur when the underlying operator $d$ has scale-invariant structure.
The log-scale framework provides the natural setting for detecting and characterizing such periodicity.
